% GENERATED FILE. Edit canonical lessons, then run npm run generate:books.
% canonical-source-hash: fnv1a64:4f3e62f28da8a0ad
% canonical-lessons: TA-C384-avaciyam, TA-C384-kattayam, TA-C384-varampu, TA-C384-kuraintapatcam, TA-C384-anaiyitu

\chapter{Musts and Limits}
\label{ch:ta-musts-and-limits}
\hlchaptermodality{384}

\begin{chapteropening}
\textbf{By the end of this chapter:} \emph{I can say what is necessary, compulsory and the limit.}

Complete the last lesson of chapter 384: I can say what is necessary, compulsory and the limit.
\end{chapteropening}

\section[\texorpdfstring{\ta{அவசியம்} (avaciyam)}{அவசியம் (avaciyam)}]{\ta{அவசியம்} (avaciyam) — necessity; necessary}

\label{lesson:TA-C384-avaciyam}



\begin{quote}
Before the new one: say the Tamil for an occasion, then the Tamil for or else.
\end{quote}

\subsection*{You'll want to know: \ta{அவசியம்}}
\textbf{\ta{அவசியம்}} — \emph{avaciyam} — \textquotedblleft{}necessity; necessary\textquotedblright{}.

\begin{quote}
\ta{தேர்வுக்கு} \ta{முன்} \ta{நன்றாகத்} \ta{தூங்குவது} \ta{அவசியம்}.
\end{quote}

\begin{quote}
\emph{tērvukku muṉ naṉṟākat tūṅkuvatu avaciyam.}
\end{quote}

\begin{quote}
It is necessary to sleep well before an exam.
\end{quote}

\ta{அவசியமான}, necessary; \ta{அவசியமில்லை}, it is not necessary.

\begin{cousinweb}
\textbf{\ta{அவசியம்}} (\emph{avaciyam}) is borrowed from Sanskrit \emph{avaśya-}, as in \emph{avaśyam}, \textquotedblleft{}necessarily, inevitably\textquotedblright{} (Monier-Williams). Tamil writes the ś as \ta{ச}.
\end{cousinweb}

\begin{grammarlens}[title={What you've built}]
One more word for saying what happens if.
\end{grammarlens}

\subsection*{Your turn}
\begin{itemize}
\raggedright
  \item \emph{Say it:} \emph{avaciyam}
  \item \emph{Say it:} \emph{avaciyam}, once more
  \item \emph{Say it:} \emph{tērvukku muṉ naṉṟākat tūṅkuvatu avaciyam.} — It is necessary to sleep well before an exam.
  \item \emph{Recall:} say the Tamil for an occasion, then the Tamil for or else, then say \emph{avaciyam} again
\end{itemize}

\begin{tcolorbox}[breakable,colback=teal!4,colframe=teal!35!black,title={Before you move on}]
What does \ta{அவசியம்} mean? (\textquotedblleft{}Necessity; necessary\textquotedblright{}.) Say it once more.
\end{tcolorbox}
\section[\texorpdfstring{\ta{கட்டாயம்} (kaṭṭāyam)}{கட்டாயம் (kaṭṭāyam)}]{\ta{கட்டாயம்} (kaṭṭāyam) — compulsion; compulsory}

\label{lesson:TA-C384-kattayam}



\begin{quote}
Before the new one: say the Tamil for or else, then the Tamil for necessity.
\end{quote}

\subsection*{You'll want to know: \ta{கட்டாயம்}}
\textbf{\ta{கட்டாயம்}} — \emph{kaṭṭāyam} — \textquotedblleft{}compulsion; compulsory\textquotedblright{}.

\begin{quote}
\ta{இந்தக்} \ta{கூட்டத்துக்கு} \ta{எல்லோரும்} \ta{கட்டாயம்} \ta{வர} \ta{வேண்டும்}.
\end{quote}

\begin{quote}
\emph{intak kūṭṭattukku ellōrum kaṭṭāyam vara vēṇṭum.}
\end{quote}

\begin{quote}
Everyone must come to this meeting.
\end{quote}

As an adverb, for sure: \ta{கட்டாயம்} \ta{வருவேன்}, I will definitely come. \ta{கட்டாயமாக}, compulsorily.

\begin{cousinweb}
\textbf{\ta{கட்டாயம்}} (\emph{kaṭṭāyam}) is native Dravidian: DEDR 1157 lists Tamil \emph{kaṭṭāyam}, \textquotedblleft{}force, compulsion, constraint; certainly\textquotedblright{}, with Kannada \emph{kaḍḍāya}, \textquotedblleft{}force\textquotedblright{}.
\end{cousinweb}

\begin{grammarlens}[title={What you've built}]
One more word for saying what happens if.
\end{grammarlens}

\subsection*{Your turn}
\begin{itemize}
\raggedright
  \item \emph{Say it:} \emph{kaṭṭāyam}
  \item \emph{Say it:} \emph{kaṭṭāyam}, once more
  \item \emph{Say it:} \emph{intak kūṭṭattukku ellōrum kaṭṭāyam vara vēṇṭum.} — Everyone must come to this meeting.
  \item \emph{Recall:} say the Tamil for or else, then the Tamil for necessity, then say \emph{kaṭṭāyam} again
\end{itemize}

\begin{tcolorbox}[breakable,colback=teal!4,colframe=teal!35!black,title={Before you move on}]
What does \ta{கட்டாயம்} mean? (\textquotedblleft{}Compulsion; compulsory\textquotedblright{}.) Say it once more.
\end{tcolorbox}
\section[\texorpdfstring{\ta{வரம்பு} (varampu)}{வரம்பு (varampu)}]{\ta{வரம்பு} (varampu) — a limit}

\label{lesson:TA-C384-varampu}



\begin{quote}
Before the new one: say the Tamil for necessity, then the Tamil for compulsion.
\end{quote}

\subsection*{You'll want to know: \ta{வரம்பு}}
\textbf{\ta{வரம்பு}} — \emph{varampu} — \textquotedblleft{}a limit\textquotedblright{}.

\begin{quote}
\ta{இந்தச்} \ta{சாலையில்} \ta{வேக} \ta{வரம்பு} \ta{நாற்பது}.
\end{quote}

\begin{quote}
\emph{intac cālaiyil vēka varampu nāṟpatu.}
\end{quote}

\begin{quote}
The speed limit on this road is forty.
\end{quote}

\ta{வயது} \ta{வரம்பு}, an age limit; \ta{வரம்பு} \ta{மீறு}, to overstep the limit.

\begin{cousinweb}
\textbf{\ta{வரம்பு}} (\emph{varampu}) is native Dravidian: DEDR 5261 lists Tamil \emph{varampu}, \textquotedblleft{}boundary, limit; the ridge of a field; rule of conduct\textquotedblright{}.
\end{cousinweb}

\begin{grammarlens}[title={What you've built}]
One more word for saying what happens if.
\end{grammarlens}

\subsection*{Your turn}
\begin{itemize}
\raggedright
  \item \emph{Say it:} \emph{varampu}
  \item \emph{Say it:} \emph{varampu}, once more
  \item \emph{Say it:} \emph{intac cālaiyil vēka varampu nāṟpatu.} — The speed limit on this road is forty.
  \item \emph{Recall:} say the Tamil for necessity, then the Tamil for compulsion, then say \emph{varampu} again
\end{itemize}

\begin{tcolorbox}[breakable,colback=teal!4,colframe=teal!35!black,title={Before you move on}]
What does \ta{வரம்பு} mean? (\textquotedblleft{}A limit\textquotedblright{}.) Say it once more.
\end{tcolorbox}
\section[\texorpdfstring{\ta{குறைந்தபட்சம்} (kuṟaintapaṭcam)}{குறைந்தபட்சம் (kuṟaintapaṭcam)}]{\ta{குறைந்தபட்சம்} (kuṟaintapaṭcam) — at least; a minimum}

\label{lesson:TA-C384-kuraintapatcam}



\begin{quote}
Before the new one: say the Tamil for compulsion, then the Tamil for a limit.
\end{quote}

\subsection*{You'll want to know: \ta{குறைந்தபட்சம்}}
\textbf{\ta{குறைந்தபட்சம்}} — \emph{kuṟaintapaṭcam} — \textquotedblleft{}at least; a minimum\textquotedblright{}.

\begin{quote}
\ta{அவன்} \ta{தினமும்} \ta{குறைந்தபட்சம்} \ta{ஒரு} \ta{மணி} \ta{நேரம்} \ta{படிக்கிறான்}.
\end{quote}

\begin{quote}
\emph{avaṉ tiṉamum kuṟaintapaṭcam oru maṇi nēram paṭikkiṟāṉ.}
\end{quote}

\begin{quote}
He studies for at least an hour every day.
\end{quote}

\ta{குறைந்த}, less + \ta{பட்சம்}, side: on the lower side.

\begin{cousinweb}
\textbf{\ta{குறைந்தபட்சம்}} (\emph{kuṟaintapaṭcam}) is built from native \emph{kuṟainta}, \textquotedblleft{}reduced\textquotedblright{}, from \emph{kuṟai} (DEDR 1851), and \emph{paṭcam}, \textquotedblleft{}side\textquotedblright{}, borrowed from Sanskrit \emph{pakṣa}, \textquotedblleft{}wing, side\textquotedblright{} (Monier-Williams).
\end{cousinweb}

\begin{grammarlens}[title={What you've built}]
One more word for saying what happens if.
\end{grammarlens}

\subsection*{Your turn}
\begin{itemize}
\raggedright
  \item \emph{Say it:} \emph{kuṟaintapaṭcam}
  \item \emph{Say it:} \emph{kuṟaintapaṭcam}, once more
  \item \emph{Say it:} \emph{avaṉ tiṉamum kuṟaintapaṭcam oru maṇi nēram paṭikkiṟāṉ.} — He studies for at least an hour every day.
  \item \emph{Recall:} say the Tamil for compulsion, then the Tamil for a limit, then say \emph{kuṟaintapaṭcam} again
\end{itemize}

\begin{tcolorbox}[breakable,colback=teal!4,colframe=teal!35!black,title={Before you move on}]
What does \ta{குறைந்தபட்சம்} mean? (\textquotedblleft{}At least; a minimum\textquotedblright{}.) Say it once more.
\end{tcolorbox}
\section[\texorpdfstring{\ta{ஆணையிடு} (āṇaiyiṭu)}{ஆணையிடு (āṇaiyiṭu)}]{\ta{ஆணையிடு} (āṇaiyiṭu) — to command; to order}

\label{lesson:TA-C384-anaiyitu}



\begin{quote}
Before the new one: say the Tamil for a limit, then the Tamil for at least.
\end{quote}

\subsection*{You'll want to know: \ta{ஆணையிடு}}
\textbf{\ta{ஆணையிடு}} — \emph{āṇaiyiṭu} — \textquotedblleft{}to command; to order\textquotedblright{}.

\begin{quote}
\ta{நீதிபதி} \ta{அவரை} \ta{விடுதலை} \ta{செய்யுமாறு} \ta{ஆணையிட்டார்}.
\end{quote}

\begin{quote}
\emph{nītipati avarai viṭutalai ceyyumāṟu āṇaiyiṭṭār.}
\end{quote}

\begin{quote}
The judge ordered him to be set free.
\end{quote}

Past (I …): \ta{ஆணையிட்டேன்}, \emph{āṇaiyiṭṭēṉ}.

The action takes -\ta{உமாறு}: \ta{செல்லுமாறு} \ta{ஆணையிடு}, to order someone to go.

\begin{cousinweb}
\textbf{\ta{ஆணையிடு}} (\emph{āṇaiyiṭu}) is built from \emph{āṇai}, \textquotedblleft{}command\textquotedblright{}, borrowed from Sanskrit \emph{ājñā} (CDIAL 1095), whose Prakrit form is \emph{āṇā}, and native \emph{iṭu}, \textquotedblleft{}to put\textquotedblright{} (DEDR 442).
\end{cousinweb}

\begin{grammarlens}[title={What you've built}]
One more word for saying what happens if.
\end{grammarlens}

\subsection*{Your turn}
\begin{itemize}
\raggedright
  \item \emph{Say it:} \emph{āṇaiyiṭu}
  \item \emph{Say it:} \emph{āṇaiyiṭu}, once more
  \item \emph{Say it:} \emph{nītipati avarai viṭutalai ceyyumāṟu āṇaiyiṭṭār.} — The judge ordered him to be set free.
  \item \emph{Recall:} say the Tamil for a limit, then the Tamil for at least, then say \emph{āṇaiyiṭu} again
\end{itemize}

\begin{tcolorbox}[breakable,colback=teal!4,colframe=teal!35!black,title={Before you move on}]
What does \ta{ஆணையிடு} mean? (\textquotedblleft{}To command; to order\textquotedblright{}.) Say it once more.
\end{tcolorbox}
